% Imporditud: tools/import_eksperimendid.py
% Allikas: Eesti füüsikaolümpiaadi eksperimendi ülesannete
%   kogu (Rael Kalda, 2023), ülesanne Ü108, lahendus L108.
\setAuthor{Eero Uustalu}
\setRound{piirkonnavoor}
\setYear{2014}
\setNumber{G E1}
\setDifficulty{3}
\setTopic{vedelike mehaanika}
\setVahendid{Joogikõrs, joonlaud, katseklaas kahte kihti jaotunud vedelikuga}
\setPunktid{10}
\setAllikas{Eksperimendi kogu Ü108, lahendus lk 83}
\setLahendusseis{toores}

\prob{Vedeliku tihedus}
Määrata kollase vedeliku tihedus kahte kihti jaotunud vedelikus. Raskema vedeliku tihedus $\rho$ = 1,0 g/cm$^3$.
\textit{Märkus:} Katses kasutatud vedelik võib määrida riideid ja paberit.

\hint

\solu
Ülesanne taandub vedelikusammaste rõhkude tasakaalule. Kõrre sisse peab tekitama teise vedelikusammaste kõrguste jaotuse kui kõrrest väljas (nagu lähtub jooniselt). Juhtude 1 ja 3 tekitamiseks peab ülumise vedeliku kõrre sisse imema, kõrre pealt sulgema ning uputama alumise otsa alumise (tihedama) vedeliku põhjani ning seejärel avama. Juhtude 2 ja 4 tekitamiseks peab kas tühja kõrre pealt sulgema ja uputama alumise otsa alumise (tihedama) vedelikuni ning seejärel avama või puhuma põhja ulatuva kõrre tühjaks ja lubama taastäituda. Katse selgitus olemas või katse läbi viidud ehk sambad tekitatud ja mõõdetud --- 2,5 p.

Märgitakse, et vajalik on eralduspinna selge eristumine ja/või oodatakse vedelike eralduspinna selginemist (1 p). Lähtudes rõhkude võrdsusest toru alumises otsas juhu 4 näitel:

p = hV A $\cdot \rho$A $\cdot$ g + hV B $\cdot \rho$B $\cdot$ g = hSA $\cdot \rho$A $\cdot$ g + hSB $\cdot \rho$B $\cdot$ g (1 p)

(hV A $-$ hSA) $\cdot \rho$A = (hSB $-$ hV B) $\cdot \rho$B

$\rho$A = hSB $-$ hV B $\cdot \rho$B (0,5 p) hV A $-$ hSA

Kõrre täiteks kasutatakse vaid ühte vedelikest (nagu juhtudel 1 ja 2 , kuna siis saab

suurima vedelikusammaste kõrguste vahe) (1 p).

Näidatakse, et täpseima tulemuse saab, kui kõrs täita vedelikuga mille kihi paksus

anumas on väiksem (kuna siis saab suurima absoluutse vedelikusammaste kõr-

guste vahe) (1 p).

Viiakse läbi vähemalt üks kontrollmõõtmine taastäidetud kõrrega (1 p). Vastus täpsusega piirides kuni $\pm_3$\% (2 p). Täpsusega $\pm_6$\% (1 punkt)

Ülejäänud katsevariandid: Variant kus kõrs üritatakse täita kummagi vedelikuga ja tekitada U-toru kokku maksiumaalselt 5 punkti. Variandid kus üritatakse kasutada kõrre sulgemist ja õhurõhku (või õhu rõhu alandamist) kokku maksimaalselt 3 punkti.

Seletuseks gümnaasiumile: viime läbi vedeliku A tiheduse mõõtemääramatuse hinnangu

$\Delta\rho$A = $\Delta$hSB + $\Delta$hV B + $\Delta$hV A + $\Delta$hSA . $\rho$A hSB $-$ hV B hV A $-$ hSA

Sellest järeldub, et vähima mõõtemääramatuse saame kui: meil on mõnes sambas puudu üks vedelikukiht; sammaste absoluutsed kõrguste erinevused on maksimaalsed võimalikest.
\probend
